\section{Gaussian limits for the full Cayley--Jacobi trace}
\label{sec:cayley-gaussian}
The fourth-moment results admit an all-order extension in sparse regimes.
The external arithmetic inputs are Chowla's cotangent independence and,
for the density statement only, Bombieri--Vinogradov and the prime number
theorem. Priority for this specialization has not been established.

Let $K$ be an odd prime, $p\equiv1\pmod K$ a prime, $H\le\Fp^*$ of order
$K$, and $A=c(H\setminus\{1\})$, where $c(h)=(h+1)/(h-1)$.
Put $n=K-1$, $d=n/2$, and
\[
 X_{K,p}(q)=n^{-1/2}\sum_{a\in A}e^{2\pi iqa/p},\qquad q\in\Fp.
\]
This is real since $A=-A$, and $X_{K,p}(0)=\sqrt n$.
Probability within a design means uniform sampling of $q\in\Fp^*$,
not sampling matrix eigenvalues.

\begin{theorem}[Pointwise sparse trace limit]
For every sequence of split prime designs with $K\to\infty$ and
$K=o(\log p)$, the law of $X_{K,p}(q)$ tends to $N(0,1)$.
For fixed distinct positive integers $m_1,\ldots,m_L$, more generally,
\[
 (X_{K,p}(m_1q),\ldots,X_{K,p}(m_Lq))\Longrightarrow N(0,I_L).
\]
\end{theorem}

\begin{theorem}[Density trace limit]
Let $Q=Q(X)\to\infty$, $\log Q=o(\log X)$, and
\[
 \mathcal D(X)=\{(K,p):K,p\text{ prime},\
 Q/2<K\le Q,\ X/2<p\le X,\ p\equiv1\pmod K\}.
\]
Let $\mu_{K,p}^{\boldsymbol m}$ be the empirical joint law in the preceding
theorem, for fixed distinct positive multipliers. For each $\varepsilon>0$,
\[
 \frac{\#\{(K,p)\in\mathcal D(X):
 d_{\rm BL}(\mu_{K,p}^{\boldsymbol m},N(0,I_L))>\varepsilon\}}
 {|\mathcal D(X)|}\longrightarrow0,
\]
where $d_{\rm BL}$ denotes the bounded-Lipschitz metric on probability measures.
\end{theorem}

These statements transfer to the full annihilator twisted-Jacobi trace through
the exact identity from the arithmetic core,
\[
 X_{K,p}(mq)=\frac K{\sqrt{K-1}}\operatorname{Tr}M_{H^\perp}(2mq).
\]
They concern parameter multiples, not $\operatorname{Tr}(M_{H^\perp}(q)^m)$.
They imply no microscopic BGS or desymmetrized spectral universality theorem.

\subsection{Weighted relations and their resultant budget}
Fix a nonempty word $\boldsymbol w=(w_1,\ldots,w_r)$ of positive integers,
and set $W=\sum_iw_i$. A tuple $\boldsymbol a\in(\mathbf F_K^*)^r$ is
balanced when
\[
 \sum_{i:a_i=b}w_i=\sum_{i:a_i=-b}w_i\qquad(1\le b\le d).
\]
Write $B_{K,\boldsymbol w}$ for its balanced count. Since
$c(\zeta^a)=-i\cot(\pi a/K)$, Chowla's independence theorem says that
$\sum_iw_ic(\zeta^{a_i})=0$ holds exactly for balanced tuples.
Wang's \emph{On a theorem of S.~Chowla}, J.~Number Theory 15 (1982), 1--4,
Theorem A, states precisely this prime-order independence input.

Let $Z_{K,\boldsymbol w}(p)$ count all ordered tuples satisfying
$\sum_iw_ic(h^{a_i})=0$ in $\Fp$, where $h$ generates $H$, and put
$\Delta_{K,\boldsymbol w}(p)=Z_{K,\boldsymbol w}(p)-B_{K,\boldsymbol w}\ge0$.
The inequality follows because every balanced relation survives reduction.
For representatives $1\le a_i<K$, use the integer polynomial
\[
 F_{\boldsymbol a,\boldsymbol w}(T)
 =\sum_{i=1}^r w_i(T^{a_i}+1)\prod_{j\ne i}(T^{a_j}-1).
\]
For a nonbalanced tuple, its resultant $R=\Res(\Phi_K,F)$ satisfies
\begin{equation}
 R\ne0,\qquad |R|\le(W2^r)^n,\qquad K^{r-1}\mid R.
 \label{eq:cayley-allorder-norm}
\end{equation}
Indeed $\|F\|_1\le W2^r$, and all $n$ conjugates of $F(\zeta)$ have this
modulus bound. Each summand of $F$ is divisible by $(T-1)^{r-1}$, so
multiplicativity of the resultant and $\Res(\Phi_K,T-1)=K$ prove divisibility.

Simultaneous scaling of the exponents by $\mathbf F_K^*$ acts freely,
preserves balance, and gives $(n^r-B_{K,\boldsymbol w})/n$ nonbalanced orbits.
Choose one representative with $a_1=1$ from each orbit.
At a split prime $p$, the $n$ roots of $\Phi_K$ lift to $\mathbf Z_p$.
Every colliding member of the orbit contributes a factor divisible by $p$
to the product formula for its representative's resultant. Hence its number
of colliding members is at most $v_p(R)$.
Since $p\ne K$, removing the factor $K^{r-1}$ and summing yields
\begin{proposition}[All-order weighted budget]
For $r\ge1$,
\begin{equation}
 \sum_{\substack{p\ {\rm prime}\\p\equiv1\pmod K}}
 \Delta_{K,\boldsymbol w}(p)\log p
 \le\frac{n^r-B_{K,\boldsymbol w}}n
 \left[n\log(W2^r)-(r-1)\log K\right]
 =O_{\boldsymbol w}(K^r).
 \label{eq:cayley-allorder-budget}
\end{equation}
Also $p>(W2^r)^n$ implies $\Delta_{K,\boldsymbol w}(p)=0$.
\end{proposition}
The zeroth moment equals one and is not included in this resultant argument.
For the scalar fourth moment, the earlier specialized coefficient $12$
is sharper; the present gain is access to every fixed weighted order.

\subsection{Balanced moments and the Gaussian comparison}
For the scalar word of length $2s$, write $B_{K,2s}$. Counting equal positive
and negative multiplicities gives
\[
 B_{K,2s}=(2s)![t^s]\left(\sum_{j\ge0}\frac{t^j}{(j!)^2}\right)^d,
 \qquad B_{K,2s+1}=0.
\]
Thus $B_{K,2}=n$, $B_{K,4}=3n(n-1)$, and
$B_{K,6}=15n^3-45n^2+40n$. The terms with $s$ different antipodal
coordinates give $(2s)!\binom ds$; the others are $O_s(d^{s-1})$.
Consequently $B_{K,2s}/n^s=(2s-1)!!+O_s(n^{-1})$.

Additive orthogonality gives the exact mixed-moment identity
\begin{equation}
 \frac1{p-1}\sum_{q\ne0}\prod_{i=1}^rX_{K,p}(w_iq)
 =\frac{p(B_{K,\boldsymbol w}+\Delta_{K,\boldsymbol w}(p))-n^r}
 {(p-1)n^{r/2}}.
 \label{eq:cayley-allorder-moment}
\end{equation}
The subtraction removes $q=0$. The two lowest cases have no excess at any
split prime, by direct $\mathbf F_p$ computations. For the word $(1)$, a
relation $c(h^a)=0$ forces $h^a=-1$, since $c(x)=(x+1)/(x-1)$ vanishes only
at $x=-1$; but $-1\notin H$, as $H$ has odd order $K$ while $-1$ has order
two. Hence $Z_{K,(1)}(p)=B_{K,(1)}=0$. For the word $(1,1)$,
$c(x)+c(y)=0$ with $x,y\in H\setminus\{1\}$ becomes, after clearing
denominators (legitimate since $x,y\ne1$ and $p$ is odd), the equation
$2xy-2=0$, i.e.\ $xy=1$; these are exactly the balanced tuples. Hence
$\Delta_{K,(1)}(p)=\Delta_{K,(1,1)}(p)=0$ for every split $p$, and
\eqref{eq:cayley-allorder-moment} gives in particular the exact mean
$-\sqrt n/(p-1)$ and the second moment $(p-n)/(p-1)$.

For independent uniform phases $\Theta_1,\ldots,\Theta_d$ on $[0,2\pi]$, put
\[
 Y_{K,j}=n^{-1/2}\sum_{b=1}^d2\cos(m_j\Theta_b).
\]
Expansion and phase integration show that their mixed moment for a word
$\boldsymbol w$ is exactly $B_{K,\boldsymbol w}/n^{r/2}$.
The vectors $(\sqrt2\cos(m_j\Theta_b))_{j=1}^L$ are bounded, independent
over $b$, have mean zero and covariance $I_L$. In a fixed mixed-moment
expansion, singleton blocks vanish. Pair partitions supply Wick's formula;
all other partitions contribute $O_{\boldsymbol w}(d^{-1/2})$ or less.
Their moments therefore tend to those of $N(0,I_L)$.
Relations among the multipliers only affect lower-order terms:
$B_{K,(1,1,2)}=n$ gives normalized mixed third moment $n^{-1/2}$.

For a fixed word, the difference between
\eqref{eq:cayley-allorder-moment} and its independent-phase moment is at most
\[
 \frac p{p-1}\frac{\Delta_{K,\boldsymbol w}(p)}{n^{r/2}}
 +\frac{2n^{r/2}}{p-1}.
\]
Under $K=o(\log p)$, the excess eventually vanishes for each fixed word,
by \eqref{eq:cayley-allorder-norm}. The remaining error tends to zero.
Gaussian moment determinacy, or the scalar moment argument for every linear
combination, proves the pointwise theorem.

\subsection{The density denominator and the law-level limit}
Bombieri--Vinogradov at $X$ and $X/2$ gives
\[
 |\mathcal D(X)|=
 (\pi(X)-\pi(X/2))
 \sum_{\substack{Q/2<K\le Q\\K\ {\rm prime}}}\frac1{K-1}
 +O_A(X/(\log X)^A).
\]
Its modulus range contains $Q=X^{o(1)}$ eventually. PNT and partial summation
give the reciprocal sum $(\log2+o(1))/\log Q$. Therefore
\[
 |\mathcal D(X)|\sim\frac{\log2}{2}\frac X{\log X\log Q}.
\]
A fixed $A=3$ already makes the relative error tend to zero, even for
arbitrarily slowly growing $Q$. No individual varying-modulus PNT is used.

By \eqref{eq:cayley-allorder-budget} and $\log p\ge\log(X/2)$,
\[
 \sum_{(K,p)\in\mathcal D(X)}
 \frac{\Delta_{K,\boldsymbol w}(p)}{(K-1)^{r/2}}
 \ll_{\boldsymbol w}
 \frac1{\log X}\sum_{\substack{Q/2<K\le Q\\K\ {\rm prime}}}K^{r/2}
 \ll_{\boldsymbol w}\frac{Q^{r/2+1}}{\log X\log Q}.
\]
After division by the design count, the average normalized excess is
$O_{\boldsymbol w}(Q^{r/2+1}/X)=o(1)$. Markov's inequality proves convergence
in probability of every mixed moment to its Gaussian value.

For completeness, this implies convergence of the random empirical laws.
From any sequence, extract a subsequence on which the first $j$ moments
are simultaneously accurate with failure probability at most $2^{-j}$.
On a product coupling, Borel--Cantelli gives almost sure convergence of all
mixed moments along this subsequence. Second moments give tightness, higher
even moments give uniform integrability of lower polynomials, and Gaussian
moment determinacy identifies every limit. Thus the bounded-Lipschitz
distance tends to zero almost surely along a further subsequence, proving
convergence in probability as claimed.

At a fixed scale $Q=X^\theta$, the same calculation proves only the
order-$r$ moment conclusion when
$0<\theta<\min(1/2,2/(r+2))$. The allowed thresholds tend to zero with $r$;
this argument does not give a fixed-polynomial-scale CLT.

\subsection{A finite obstruction to the fourth-moment shortcut}
Exact integer convolution for $(K,p)=(11,199)$ gives
\[
 (Z_3,B_3)=(6,0),\qquad (Z_4,B_4)=(270,270),\qquad
 (Z_6,B_6)=(12250,10900).
\]
Thus the exact fourth-order baseline does not force the higher-order
classification. The all-order budgets are needed.
The maintained verifier checks 32 finite word/order designs, representing
6440 ordered tuples and 1242 scaling orbits, with 13293 orbit-valuation
checks (run log
\path{evidence/executions/20261006/corpus-revision/final/checks/cayley_higher_moments.log}). These are implementation regressions, not proofs of the external
arithmetic inputs. The full statements, proof details, source locators,
and limits of the BGS interpretation are in
\path{mathematics/cayley_trace_gaussian_limit.md}.

\subsection{A quantitative scalar refinement}
The scalar limit admits a rate by keeping all constants uniform in the
number of moments. Write $\mu_{K,p}$ for the empirical scalar law with
$q\ne0$, $\gamma=N(0,1)$, and
$D_{K,p}=\sup_x|\mu_{K,p}((-\infty,x])-\gamma((-\infty,x])|$.

\begin{theorem}[Pointwise Kolmogorov bound]
There is an absolute constant $C$ such that, for every odd prime $K$ and
split prime $p$,
\[
 D_{K,p}\le C\left((K-1)^{-1/2}
                    +\sqrt{\frac{K-1}{\log p}}\right).
\]
\end{theorem}

\begin{theorem}[Density Kolmogorov bound]
For sufficiently large $X,Q$ with $Q\le X^{1/4}$, all but a proportion
at most $C(\log X)X^{-1/4}$ of the designs in $\mathcal D(X)$ satisfy
\[
 D_{K,p}\le C\left(Q^{-1/2}
                   +\sqrt{\frac{\log Q}{\log X}}\right).
\]
The average distance is bounded by this right side plus
$C(\log X)X^{-1/4}$. Constants are absolute and not optimized.
\end{theorem}
For fixed $Q=X^\theta$ this only gives an $O(\sqrt\theta)$ discrepancy;
it does not establish a fixed-power CLT.

\paragraph{Finite moments and smoothing.}
Let $\mu^{\rm all}_{K,p}$ include $q=0$, and let $\nu_K$ be the scalar
independent-phase law above. If $\Delta_{K,r}(p)=0$ for $1\le r\le2s$,
then these laws have exactly the same first $2s$ moments.
The coefficientwise inequality
\[
 \sum_{j\ge0}\frac{z^{2j}}{n^j(j!)^2}\preceq\exp(z^2/n)
\]
is preserved by taking the $d$-th power. Thus
$\int y^{2s}\,d\nu_K(y)\le(2s)!/(2^ss!)$.
Taylor expansion of $e^{ity}$ through degree $2s-1$ gives
\[
 |\widehat\mu^{\rm all}_{K,p}(t)-\widehat\nu_K(t)|
 \le\frac{2|t|^{2s}}{2^ss!}.
\]
For $T^2\le s/(2e)$ this implies
\[
 \int_0^T\frac{|\widehat\mu^{\rm all}_{K,p}(t)-\widehat\nu_K(t)|}{t}\,dt
 \le\frac{T^{2s}}{s\,2^ss!}\le\frac{4^{-s}}s.
\]

For $|t|\le\sqrt n/4$, put $v=t^2/n\le1/16$.
The cosine expansion gives
$1-v\le J_0(2t/\sqrt n)\le1-v+v^2/4$, so this quantity is positive and
\[
 |\log J_0(2t/\sqrt n)+v|\le\tfrac8{15}v^2.
\]
Both $d\log J_0(2t/\sqrt n)$ and $-t^2/2$ are at most $-t^2/3$.
Since $\widehat\nu_K(t)=J_0(2t/\sqrt n)^d$, the mean-value theorem gives
\[
 |\widehat\nu_K(t)-e^{-t^2/2}|
 \le\frac4{15n}t^4e^{-t^2/3},\qquad
 \int_0^T\frac{|\widehat\nu_K(t)-e^{-t^2/2}|}{t}\,dt\le\frac6{5n}.
\]
The logarithmic bound follows from
$\log(1-v)\ge-v-v^2/[2(1-v)]$ and $\log(1+u)\le u$.

Apply the classical Esseen smoothing inequality against the normal target
with $T=\min\{\sqrt{s/(2e)},\sqrt n/4\}$. This inequality permits an atomic
first law; only the target needs a Lipschitz CDF. An exact statement is
B.~Borda, arXiv:2005.04925v2, Theorem A, p.~2.
The smoothing remainder $O(1/T)$ and the preceding integrals imply
\[
 d_{\rm K}(\mu^{\rm all}_{K,p},\gamma)\le C(s^{-1/2}+n^{-1/2}).
\]
Removing zero changes the law by at most $1/p$ in Kolmogorov distance, because
$\mu^{\rm all}_{K,p}=p^{-1}\delta_{\sqrt n}+(1-p^{-1})\mu_{K,p}$. Hence
\begin{equation}
 \big[\Delta_{K,r}(p)=0\ (1\le r\le2s)\big]
 \ \Longrightarrow\
 D_{K,p}\le C(s^{-1/2}+n^{-1/2})+1/p.
 \label{eq:cayley-finite-moment-rate}
\end{equation}

\paragraph{Uniform arithmetic cutoffs.}
For $r\le2s$, $r2^r\le2s\,2^{2s}\le2^{4s}$.
Thus $p>2^{4ns}$ excludes every nonbalanced scalar relation through order
$2s$. Taking $s=\lfloor\log p/(8n\log2)\rfloor$ when it is positive
ensures the strict cutoff and gives the pointwise theorem by
\eqref{eq:cayley-finite-moment-rate}. If the floor is zero, the right side
of the asserted bound is bounded below by an absolute positive constant,
and $D_{K,p}\le1$ suffices.

For density, the all-order budget and $\log(r2^r)\le2r$ imply uniformly
for every integer $s\ge1$ that
\[
 \sum_{r=1}^{2s}\sum_{p\equiv1\pmod K}\Delta_{K,r}(p)\log p
 \le\sum_{r=1}^{2s}2r n^r\le8s n^{2s}.
\]
A design with any excess through order $2s$ consumes at least $\log(X/2)$,
by integrality. Bound $n^{2s}\le Q^{2s}$ before using the prime count
$O(Q/\log Q)$. Dividing by the same uniform Bombieri--Vinogradov
denominator shows that the exceptional fraction is at most
$CsQ^{2s+1}/X$, with a constant independent of $s$.
For $s=\lfloor\log X/(4\log Q)\rfloor$ and $Q\le X^{1/4}$,
\[
 s\ge1,\qquad Q^{2s+1}\le X^{3/4},\qquad
 s^{-1/2}\le C\sqrt{\log Q/\log X}.
\]
The exceptional fraction is therefore at most $C(\log X)X^{-1/4}$.
On its complement, use \eqref{eq:cayley-finite-moment-rate},
$n\ge Q/3$ and $p>X/2$. This proves the density rate and, by bounding
distance by one on the exceptional set, its averaged version.
No dimension-uniform multivariate rate or microscopic spectral law is inferred.
